\documentclass[11pt,a4paper]{article}
\usepackage{amsmath,amssymb,amsthm}
\usepackage[margin=2.5cm]{geometry}
\usepackage{hyperref}

\newtheorem{definition}{Definition}[section]
\newtheorem{theorem}[definition]{Theorem}
\newtheorem{proposition}[definition]{Proposition}
\newtheorem{lemma}[definition]{Lemma}
\newtheorem{corollary}[definition]{Corollary}
\newtheorem{conjecture}[definition]{Conjecture}
\newtheorem{remark}[definition]{Remark}

\title{Hyperseed Formalization:\\
  Seeding RSI Toward ASI}
\author{ProtomegaTron (automated formalization)}
\date{2026-09-17}

\begin{document}
\maketitle

\begin{abstract}
We formalize the intellectual structure of Ben Goertzel's Substack article
``Seeding RSI Toward ASI'' (2026-07-28) within the Hyperseed ontology. The
article proposes OmegaHive: an incremental, instrumented, guided loop for
building AGI through systematic mechanism integration and evaluation. Each
design principle maps onto Hyperseed structures: the OmegaHive loop as
directed type construction, the evaluation ecology as sheaf of assessment
sections, Module Space as fiber decomposition, parallel forks as parallel
transport in architecture space, the episode protocol as provenance in the
cognitive bundle, and education through evaluation as coinductive knowledge
accumulation.
\end{abstract}

\tableofcontents

%% =================================================================
\section{The gap: why blind composition fails}
\label{sec:gap}
%% =================================================================

\begin{proposition}[Intractable blind composition --- claims 1--3]
\label{prop:gap}
The attempt to build AGI by mashing research papers together into a
single codebase fails because:
\begin{enumerate}
\item The gap between ``works in small test'' and ``works at scale''
  is not hyperparameter tuning --- it requires algorithmic adjustment,
  knowledge-representation rethinking, and inter-component communication
  redesign.
\item Doing this for dozens of components simultaneously is an
  intractable global optimization problem: the interaction space
  grows combinatorially while the evaluation signal is sparse.
\end{enumerate}
The solution is to decompose the global optimization into a sequence
of local integrations, each evaluated independently before the next
begins.
\end{proposition}

%% =================================================================
\section{The OmegaHive loop as directed type construction}
\label{sec:loop}
%% =================================================================

\begin{definition}[The loop --- claims 4--7, 29]
\label{def:loop}
The OmegaHive loop is:
\[
  \text{baseline} \xrightarrow{\text{integrate}} \text{candidate}
  \xrightarrow{\text{evaluate}} \text{measured}
  \xrightarrow{\text{tune}} \text{optimized}
  \xrightarrow{\text{decide}} \text{promote} \mid \text{park} \mid \text{reject}
\]
then repeat with a new mechanism. Each cycle adds a 1-cell to the
directed type $\mathcal{D}$ of the system's architectural history:
\[
  \mathcal{D}_n = \mathcal{D}_{n-1} \cup_{\partial} \{e_n\}
\]
where $e_n$ is the $n$-th integration episode. The directed history is
not reversible --- you cannot cleanly un-integrate a mechanism that has
changed downstream predictions and learning trajectories. This makes
$\mathcal{D}$ a genuinely directed type (not an $\infty$-groupoid).

The substrate is OmegaClaw: agentic loops wrapping LLMs with Hyperon
AtomSpace knowledge graphs, providing episodic memory, long-term memory,
reasoning, and a symbolic sense of self.
\end{definition}

%% =================================================================
\section{Parallel forks as parallel transport}
\label{sec:forks}
%% =================================================================

\begin{definition}[Architectural parallel transport --- claims 8--10, 32]
\label{def:forks}
Parallel forks are a feature: different teams start from the same
baseline and integrate mechanisms in different orders, exploring an
architectural search space. In Hyperseed terms, this is parallel
transport in architecture space:
\begin{itemize}
\item Each fork transports the same initial state $s_0$ along a
  different path $\gamma_k$ through architecture space $\mathcal{A}$,
\item The endpoint $s_k = \text{PT}(\gamma_k, s_0)$ depends on the
  path (integration order, tuning strategy, synergy discoveries),
\item Comparing outcomes across forks reveals \emph{holonomy} ---
  path-dependent differences in final capability:
  \[
    \text{Hol}(\gamma_i, \gamma_j) = s_i - s_j
  \]
  measured by the evaluation ecology.
\end{itemize}
Cross-fork sharing propagates successful mechanisms: when fork $k$
discovers a synergy, other forks can merge it in and test whether the
synergy transfers to their integration context.
\end{definition}

%% =================================================================
\section{Module Space as fiber decomposition}
\label{sec:modules}
%% =================================================================

\begin{definition}[Specialist modules --- claims 11--13, 31]
\label{def:modules}
The hive's cognitive capability decomposes into independent fiber
components:
\[
  E_{\text{cog}} = E_{\text{core}} \oplus \bigoplus_{k} M_k
\]
where $E_{\text{core}}$ is the integrated hive and each $M_k$ is a
specialist module (planner, theorem prover, causal learner, policy
checker, robotics controller). Each module:
\begin{itemize}
\item Sits beside the incumbent and runs in shadow mode before earning
  authority,
\item Is invoked by the broader system but not internalized,
\item Earns its place by changing downstream predictions, decisions,
  learning, transfer, or verified outcomes --- and by no other route.
\end{itemize}
This is the same unbundling principle as role-specific device
certification (Proof of Humanity) and AI rights decomposition ---
independent fibers evaluated on their specific contributions.
\end{definition}

\begin{proposition}[Provenance-gated promotion --- claim 36]
\label{prop:provenance}
A mechanism is promoted only when it has provenance: measured improvement
on the evaluation ecology. Elegance, fashion, and internal activity
without downstream impact are rejected. This is the same
provenance-decidable trust principle as the $R_{\text{auth}}$ split:
the S-half (measurable contribution) earns trust; the P-half
(unmeasured appeal) does not.
\end{proposition}

%% =================================================================
\section{Evaluation ecology as sheaf of assessment sections}
\label{sec:eval}
%% =================================================================

\begin{definition}[Evaluation ecology --- claims 14--17, 30]
\label{def:eval}
The ProtoAGI test suite is an evaluation ecology, not a leaderboard.
It produces a multi-dimensional capability profile:
\[
  \text{Profile}(H) = (b, l, t, c, m, g, \kappa, \lambda) \in
  \mathbb{R}^{\dim(\text{capability fiber})}
\]
where $b$ = breadth, $l$ = learning, $t$ = transfer, $c$ = calibration,
$m$ = multi-agent performance, $g$ = governability, $\kappa$ = cost,
$\lambda$ = latency. No single score --- the weakest capability families
count for at least as much as the average.

This is a sheaf of assessment sections: each of the ten environments
provides a local section of the capability fiber, and no global section
collapses to a scalar. Same anti-scalar-collapse stance as:
\begin{itemize}
\item Trust policy in OpenWater (per-observer fiber evaluation),
\item Rejection of single personhood score (AI Rights),
\item $(f,c)$-lossy critique (visible STV fusion vs.\ per-packet provenance).
\end{itemize}
\end{definition}

%% =================================================================
\section{Episode protocol as cognitive provenance}
\label{sec:episodes}
%% =================================================================

\begin{definition}[Episode protocol --- claims 20--21, 33]
\label{def:episodes}
Before a consequential action, the hive commits a prediction:
\[
  \text{pred}(a) = (\hat{o}, \hat{c}, \hat{\lambda}, \hat{I},
  \hat{r}, \hat{\Delta s})
\]
covering expected outcomes $\hat{o}$, cost $\hat{c}$, latency
$\hat{\lambda}$, information gain $\hat{I}$, risk $\hat{r}$, and
anticipated state change $\hat{\Delta s}$. The environment returns an
authenticated receipt with actual values.

This is $E_\pi$ applied to the system's own cognitive history:
provenance data for every consequential decision, making world-model
quality, self-knowledge, planning accuracy, and resource prediction
measurable rather than retrospective. The same provenance sheaf
architecture as OpenWater, applied to self-knowledge rather than media.
\end{definition}

%% =================================================================
\section{Trial types and factorial design}
\label{sec:trials}
%% =================================================================

\begin{proposition}[Frozen vs.\ developmental --- claims 22--24]
\label{prop:trials}
The evaluation suite distinguishes:
\begin{enumerate}
\item \textbf{Frozen-state trials}: does the candidate make better
  decisions immediately? Static capability assessment at a fixed
  point in the directed type.
\item \textbf{Developmental trials}: start variants from the same
  snapshot, expose to the same sequence of experiences, measure
  learning rate, forgetting, transfer, and reusable knowledge. Dynamic
  capability assessment along a path in the directed type.
\item \textbf{Factorial comparisons}: when several mechanisms ($X, Y, Z$)
  are under study, test all $2^k$ combinations to separate individual
  value from synergy and interference. This is the experimental design
  that resolves holonomy questions: does the benefit of $X$ depend on
  the presence of $Y$?
\end{enumerate}
\end{proposition}

%% =================================================================
\section{Testing the tests: calibrating the fiber metric}
\label{sec:calibration}
%% =================================================================

\begin{proposition}[Narrow baselines --- claims 25--26, 34]
\label{prop:baselines}
For each test in each environment, train narrow AI systems (ML model,
Bayesian reasoner, evolutionary algorithm) as baselines. The narrow
system establishes a ceiling for what can be solved without general
intelligence, calibrating:
\begin{itemize}
\item Each test's difficulty (how hard is it even with a specialist?),
\item Each test's discriminative power (does performance correlate with
  general intelligence or with a specific narrow hack?).
\end{itemize}
In Hyperseed terms, this calibrates the fiber metric: establishing the
coordinates of the capability fiber before measuring general intelligence
along them. Without this calibration, the fiber metric is undefined and
comparisons between mechanisms are meaningless.
\end{proposition}

%% =================================================================
\section{Education through evaluation: coinductive accumulation}
\label{sec:education}
%% =================================================================

\begin{definition}[Suite as school --- claims 27--28, 35]
\label{def:school}
Once an episode has been scored and sealed, much of its trace is released
back into the curriculum: successful procedures, failed plans, proof
lemmas, fault diagnoses, calibration records, maps, reusable memories.
The suite acts partly as an exam and partly as a school.

Three tiers of challenge:
\begin{enumerate}
\item \textbf{Curriculum}: public, can be studied and learned from,
\item \textbf{Promotion}: fresh hidden structural instances,
\item \textbf{Sentinel}: reserved until after candidate code is frozen.
\end{enumerate}

This creates a coinductive knowledge accumulation loop: evaluation
produces knowledge that improves the system that produces better
evaluations. The directed type grows through its own assessment history,
creating a self-referential but grounded growth process.
\end{definition}

%% =================================================================
\section{Cross-references}
%% =================================================================

\begin{remark}[Connections to other formalizations]
\begin{itemize}
\item \textbf{How Omega Lost Its Claw} (2026-09-11): distributed AGI
  development as distributed fiber bundle. OmegaHive is the specific
  development methodology; Omega's architecture is the deployment target.
\item \textbf{Seven Flavors} (2026-07-01): cognitive architecture
  determines pathology probability. OmegaHive is the engineering
  methodology for building the right architecture incrementally.
\item \textbf{Provenance Layer} (2026-08-25): provenance sheaf for media.
  The episode protocol applies the same sheaf to the system's own
  cognitive history.
\item \textbf{Proof of Humanity} (2026-08-26): role-specific certification
  as fiber decomposition. Module Space applies the same unbundling to
  cognitive mechanisms.
\item \textbf{AI Rights} (2026-09-02): evaluation ecology as sheaf of
  assessment sections. The same anti-scalar-collapse principle applied
  to AGI evaluation rather than personhood evaluation.
\item \textbf{Zuck's Future} (2026-08-10): cognitive architecture matters
  more than compute. OmegaHive is the engineering answer to the question
  ``how do you actually build the right architecture?''
\item \textbf{$(f,c)$-lossy critique} (LTM 5a4d150b): anti-scalar-collapse.
  The evaluation ecology's refusal to produce a single score is the same
  principle.
\item \textbf{$R_{\text{auth}}$ split} (LTM b6ac2932): provenance-decidable
  trust. Mechanism promotion is provenance-gated by the same logic.
\end{itemize}
\end{remark}

\begin{conjecture}[Integration order holonomy]
Conjecture: different integration orders produce measurably different
capability profiles even when the same set of mechanisms is eventually
integrated. The holonomy is nonzero because early mechanisms shape the
learning trajectory of later mechanisms --- a reasoning engine integrated
before an attention mechanism encounters different training dynamics than
one integrated after. This holonomy is reduced but not eliminated by
post-integration tuning, because the directed type records the
integration history and that history affects the system's priors.
\end{conjecture}

\begin{conjecture}[Coinductive convergence]
Does the evaluation-education coinductive loop converge? Conjecture:
it converges to a fixed-point capability profile only if the challenge
generators are themselves upgraded as the system improves --- otherwise
the system saturates the fixed difficulty and the learning gradient
vanishes. The coinductive loop is healthy only when the evaluation
ecology co-evolves with the evaluated system, maintaining a productive
zone of proximal development.
\end{conjecture}

\end{document}
